\documentclass[11pt]{article}

% ── Page layout (arXiv-style single-column preprint) ──────────────────────────
\usepackage[letterpaper,margin=1in]{geometry}
\usepackage[T1]{fontenc}
\usepackage[utf8]{inputenc}
\usepackage{lmodern}

% ── Math & tables ─────────────────────────────────────────────────────────────
\usepackage{amsmath,amssymb,amsfonts}
\usepackage{booktabs}
\usepackage{multirow}
\usepackage{array}
\usepackage{graphicx}
\usepackage{xcolor}
\usepackage{pifont}

% ── Typography, links, references ─────────────────────────────────────────────
\usepackage{microtype}
\usepackage{url}
\usepackage{cite}
\usepackage[hidelinks,breaklinks=true]{hyperref}
\usepackage{fancyhdr}

\newcommand{\cmark}{\ding{51}}
\newcommand{\xmark}{\ding{55}}
\newcolumntype{L}[1]{>{\raggedright\arraybackslash}p{#1}}
\newcolumntype{C}[1]{>{\centering\arraybackslash}p{#1}}

% ── Preprint running header / footer ──────────────────────────────────────────
\pagestyle{fancy}
\fancyhf{}
\renewcommand{\headrulewidth}{0.4pt}
\fancyhead[L]{\footnotesize\itshape A Reproducible Strong Baseline for Vehicle Re-ID on VeRi-776}
\fancyhead[R]{\footnotesize\itshape Preprint}
\fancyfoot[C]{\thepage}

% ── Title & Authors ───────────────────────────────────────────────────────────
\title{\bfseries 93.3\% mAP on VeRi-776: Revisiting ViT Training Recipes and\\
       Score-Level Fusion for Vehicle Re-Identification%
\thanks{Preprint. \copyright~2026 The Authors. This work is licensed under a
  Creative Commons Attribution 4.0 International (CC BY 4.0) License. Source
  code, configuration files, and evaluation notebooks are available at
  \url{https://github.com/MRKDaGods/gp}.}}

\author{
  Abdelrahman Hamdi$^{1}$ \qquad Ali Ashraf Ali$^{1}$ \qquad
  Mohamed Amar Ali$^{1,\ast}$ \\[2pt]
  Seif Tamer Heakal$^{1}$ \qquad Ahmed Hamdi$^{1}$ \\[6pt]
  \normalsize $^{1}$Department of Computer Engineering, Faculty of Engineering,
    Cairo University, Giza, Egypt \\[3pt]
  \small \texttt{abdelrahman.ibrahim03@eng-st.cu.edu.eg} \quad
    \texttt{ali.abdelraouf03@eng-st.cu.edu.eg} \\[1pt]
  \small \texttt{mohamed.awaad02@eng-st.cu.edu.eg} \quad
    \texttt{seif.mohamed05@eng-st.cu.edu.eg} \\[1pt]
  \small \texttt{ahamdy@eng.cu.edu.eg} \\[3pt]
  \normalsize $^{\ast}$\,Corresponding author
}

\date{June 2026}

\begin{document}
\maketitle
\thispagestyle{fancy}

% ── Abstract ──────────────────────────────────────────────────────────────────
\begin{abstract}
Vehicle re-identification (Re-ID) in urban surveillance is a fundamental
challenge in intelligent transportation systems, requiring robust cross-camera
matching under viewpoint variation, illumination shifts, and high inter-class
visual similarity. While transformer-based methods have substantially advanced
the field, their training recipes remain systematically under-optimized: the
\emph{de facto} ViT-B baseline, TransReID~\cite{he2021transreid}, reports
82.3\% mAP on VeRi-776, and subsequent improvements have largely been driven
by architectural novelty rather than recipe refinement.

This paper presents two complementary contributions. First, a systematically
improved training recipe for the ViT-Base/16 backbone---retaining the TransReID
SIE and JPM modules and initialized from CLIP pretrained
weights~\cite{radford2021clip}---combined with a tuned metric-learning
objective and a three-step post-processing pipeline (AQE + $k$-reciprocal
re-ranking~\cite{zhong2017rerank} + flip TTA) achieves \textbf{89.97\% mAP
and 97.79\% Rank-1} at its best-mAP operating point on VeRi-776, with a
separate best-Rank-1 operating point of \textbf{85.14\% mAP and 98.33\%
Rank-1}. This is the best reported mAP for any pure ViT-B model on this
benchmark, surpassing TransReID by $+7.67$\,pp mAP.

Second, a training-free score-level fusion of this ViT-B stream with the
CLIP-based Semantic Enhancement Network
(CLIP-SENet~\cite{lu2025clipsenet}) achieves \textbf{93.32\% mAP and
98.45\% Rank-1}. This establishes a new state of the art in mAP on VeRi-776, while
maintaining competitive Rank-1 accuracy relative to the previous best published result
(CLIP-SENet standalone: 92.9\%~mAP and 98.7\% Rank-1~\cite{lu2025clipsenet}). A comprehensive ablation
study traces the full performance trajectory from 82.3\% to 93.3\% mAP.
Checkpoints, evaluation notebooks, and configuration files are publicly
released to serve as a reproducible strong baseline for the vehicle ReID
community.
\end{abstract}

\vspace{0.6em}
\noindent\textbf{Keywords:}\;
Vehicle re-identification, Vision Transformer, CLIP, score-level fusion,
VeRi-776, average query expansion, $k$-reciprocal re-ranking, supervised
contrastive learning, intelligent transportation systems.

% =============================================================================
\section{Introduction}
\label{sec:intro}
% =============================================================================

\subsection{Motivation}
Vehicle re-identification is a core enabling technology for intelligent
transportation systems (ITS), supporting applications including urban traffic
surveillance, law enforcement, smart city infrastructure, and autonomous fleet
management~\cite{liu2016veri,liu2018provid}. The task requires matching the
same physical vehicle across multiple spatially distributed cameras whose
fields of view do not overlap, under conditions of substantial appearance
variation: the same vehicle may be observed from different angles, under
different illumination conditions, at varying resolutions, and after arbitrary
time intervals~\cite{liu2016veri}.

The VeRi-776 dataset~\cite{liu2016veri,liu2018provid}, introduced by Liu
\emph{et al.}\ at ICME 2016 and expanded through the PROVID framework
published in \emph{IEEE Transactions on Multimedia} in 2018, has become the
standard benchmark for vehicle ReID evaluation. It contains over 50\,000
images of 776 vehicles captured by 20 cameras covering a
$1.0\,\text{km}^{2}$ urban area over 24 hours, with each vehicle appearing in
2 to 18 cameras under varied viewpoints, illuminations, and
occlusions~\cite{liu2016veri,liu2018provid}. The standard evaluation protocol
employs 37\,781 training images of 576 identities, 11\,579 gallery images, and
1\,678 query images across 200 test identities~\cite{liu2018provid}.
Performance is measured by mean Average Precision (mAP) and Rank-1
Cumulative Matching Characteristic (CMC).

\subsection{The ViT-B Training Recipe Gap}
The introduction of TransReID~\cite{he2021transreid} marked a turning point in
vehicle ReID by applying a pure Vision Transformer
(ViT)~\cite{dosovitskiy2021vit} backbone to the task. TransReID augments the
standard ViT-B/16 architecture with two purpose-built modules: the Side
Information Embedding (SIE), which encodes camera index and viewpoint as
learnable token offsets to reduce domain bias, and the Jigsaw Patch Module
(JPM), which shuffles patch tokens in a parallel branch to force locally
robust feature learning~\cite{he2021transreid}. With DeiT-B/16 backbone and
SGD (momentum=0.9, LR=0.008, cosine decay), TransReID reports 82.3\% mAP and
97.1\% Rank-1 on VeRi-776~\cite{he2021transreid}.

Despite broad adoption of this architecture as a baseline, subsequent work has
primarily focused on augmenting it with new architectural
components---semantic alignment via language
models~\cite{li2023clipreid,lu2025clipsenet}, graph-based spatial
reasoning~\cite{shen2022hpgn}, or multi-branch CNN
structures~\cite{luo2023mbr4b}---rather than systematically revisiting what
the ViT-B training recipe itself leaves on the table. No prior work reports
above 90\% mAP on VeRi-776 with a pure ViT-B backbone. This paper challenges
that assumption directly.

\subsection{Contributions}
The contributions of this work are as follows:
\begin{itemize}
  \item \textbf{C1---Best Pure ViT-B Result:} An improved training recipe
    combining CLIP pretraining initialization, a tuned metric-learning
    objective, an optimized augmentation pipeline, and AQE + $k$-reciprocal
    re-ranking~\cite{zhong2017rerank} achieves \textbf{89.97\%
    mAP / 97.79\% Rank-1} at the best-mAP operating point on VeRi-776---the
    highest reported mAP for any pure ViT-B model, surpassing CLIP-ReID (84.5\%
    mAP~\cite{li2023clipreid}) by $\mathbf{+5.47}$\,\textbf{pp mAP} and
    TransReID~\cite{he2021transreid} by $\mathbf{+7.67}$\,\textbf{pp mAP}.
  \item \textbf{C2---New State of the Art in mAP:} Score-level fusion with
    CLIP-SENet~\cite{lu2025clipsenet} achieves \textbf{93.32\% mAP /
    98.45\% Rank-1}, surpassing the previous best published mAP
    (92.9\%~\cite{lu2025clipsenet}) by $\mathbf{+0.42}$\,\textbf{pp}.
    Rank-1 (98.45\%) is competitive with but does not exceed CLIP-SENet
    standalone Rank-1 (98.7\%).
  \item \textbf{C3---Full Ablation Path:} A comprehensive ablation study
    traces the complete performance trajectory from the 82.3\% TransReID
    baseline to 93.3\% mAP, attributing each gain to a specific recipe
    component or post-processing step.
  \item \textbf{C4---Reproducible Public Release:} All checkpoints (SHA-256
    pinned), evaluation notebooks, configuration files, and the
    post-processing pipeline are publicly released. Every reported number is
    the deterministic evaluation of a frozen, hash-pinned checkpoint and is
    thus reproducible \emph{exactly} at any time
    (Section~\ref{sec:reproducibility}).
\end{itemize}

\subsection{Paper Organization}
Section~\ref{sec:related} reviews related work across CNN-era methods,
transformer-based methods, and post-processing techniques.
Section~\ref{sec:method} describes the full methodology including both streams
and the fusion strategy. Section~\ref{sec:setup} details the experimental
setup. Section~\ref{sec:results} presents quantitative results and ablations.
Section~\ref{sec:discussion} discusses findings and limitations.
Section~\ref{sec:conclusion} concludes.

% =============================================================================
\section{Related Work}
\label{sec:related}
% =============================================================================

\subsection{CNN-Era Vehicle Re-Identification}
The VeRi-776 benchmark was introduced by Liu et al.~\cite{liu2016veri} and
extended with the PROVID framework~\cite{liu2018provid}, which combined visual
appearance, license plate recognition, camera topology, and spatiotemporal cues,
achieving 53.42\% mAP and 81.56\% Rank-1.\footnote{The PROVID (53.42/81.56)
and VAMI (50.13/77.03) figures are the values reported in the respective
papers and reproduced throughout the survey
literature~\cite{ye2021reidsurvey,li2024vehiclereidsurvey}.}
A broad survey of deep learning approaches to person and vehicle ReID is
provided by Ye et al.~\cite{ye2021reidsurvey}, while Li et al.~\cite{li2024vehiclereidsurvey}
survey vehicle-specific advances including part-level modelling, multi-source
training, and cross-domain adaptation.

VAMI~\cite{zhou2018vami} (CVPR 2018) synthesised target viewpoints via
adversarial training, reaching 50.13\% mAP and 77.03\% Rank-1.
PVEN~\cite{meng2020pven} (CVPR 2020) introduced parsing-based view-aware
embeddings via semantic segmentation for part-level feature extraction,
achieving 79.5\% mAP and 95.6\% Rank-1.
CAL~\cite{rao2021cal} (ICCV 2021) proposed counterfactual attention learning
to suppress spurious background cues, reaching 74.3\% mAP and 95.4\% Rank-1.
VehicleNet~\cite{zheng2020vehiclenet} (IEEE TMM 2020) demonstrated multi-source
data aggregation across VeRi-776, VehicleID, and synthetic data, achieving
83.41\% mAP and 96.78\% Rank-1.
HPGN~\cite{shen2022hpgn} (IEEE T-ITS 2022) modelled spatial significance
across image scales via a pyramid graph network, achieving 80.18\% mAP and
96.72\% Rank-1.
A Strong Baseline~\cite{cybercore2021strong} (CVPR-W AICity 2021) combined
ResNet101-IBN-a with multi-head attention for multi-domain training, reaching
87.1\% mAP and establishing a competitive ResNet-family ceiling prior to
transformer approaches. Khorramshahi et al.~\cite{khorramshahi2023selfsupervision}
later revisited scalable vehicle ReID through self-supervised pretraining across
ResNet, ViT, Swin, and ConvNeXt backbones in the AI City setting.

\subsection{Transformer-Based Vehicle Re-Identification}
TransReID~\cite{he2021transreid} (ICCV 2021) was the first pure Vision
Transformer framework for object re-identification, introducing the SIE and
JPM modules and achieving 82.3\% mAP and 97.1\% Rank-1 on VeRi-776 with SGD
optimisation. Despite its architectural contributions, the TransReID recipe
leaves substantial headroom: as we demonstrate, switching the initialisation
from ImageNet-21k to CLIP~\cite{radford2021clip} and the optimiser from SGD to
AdamW~\cite{loshchilov2019adamw} alone yields a substantial gain without any
architectural change.

CLIP-ReID~\cite{li2023clipreid} (AAAI 2023) introduced a two-stage
prompt-tuning procedure that exploits CLIP's vision-language representations
without requiring text annotations, achieving 84.5\% mAP and 97.3\% Rank-1 on
VeRi-776~\cite{opencodepapers2024veri}. CLIP-ReID uses the same ViT-B/16 CLIP
backbone family as the present work; our recipe surpasses it by $+5.47$\,pp
mAP purely through recipe differences (loss, augmentation, post-processing),
not architectural changes.
Recent transformer variants have broadened this line beyond single-domain
VeRi-style retrieval: MSFFT~\cite{holla2024msfft} fuses multi-scale transformer
features for cross-platform vehicle ReID, while MiKeCoCo~\cite{kuang2025mikecoco}
uses CLIP-guided prompt knowledge for generalizable vehicle ReID.

RPTM~\cite{ghosh2023rptm} (WACV 2023) proposed relation-preserving triplet
mining to stabilise triplet loss training across hard and semi-hard negatives,
achieving 88.0\% mAP and 97.3\% Rank-1.
MBR4B-LAI~\cite{luo2023mbr4b} (ITSC 2023) combined four convolutional branches
with location-aware identity embeddings and camera/pose metadata, reaching
92.09\% mAP and 98.03\% Rank-1 after re-ranking---a competitive CNN ensemble
that our fusion method surpasses.
FastReID~\cite{wang2020fastreid} provided a unified, production-quality baseline
toolbox widely used for ablation and comparison across vehicle and person ReID
benchmarks.

CLIP-SENet~\cite{lu2025clipsenet} (arXiv 2025) is a single-stream approach
that combines ResNet101-IBN-a local appearance features with 512-dimensional
semantic embeddings from a TinyCLIP Semantic Extraction Module (SEM) and an
Adaptive Fine-Grained Enhancement Module (AFEM), achieving 92.9\% mAP and
98.7\% Rank-1 on VeRi-776. This is the closest prior result to ours; the
present work takes the pre-trained CLIP-SENet v6 checkpoint as a frozen second
stream in a training-free score-level fusion, adding $+0.42$\,pp mAP above the
CLIP-SENet standalone result.

\subsection{Re-ranking and Post-processing}
The $k$-reciprocal encoding method by Zhong et al.~\cite{zhong2017rerank}
(CVPR 2017) leverages the observation that if a gallery image lies in the
$k$-reciprocal nearest neighbours of a probe, it is more likely to be a true
match. A Jaccard-distance-based score is blended with the original cosine
distance, improving retrieval precision without requiring labelled data.
Zhang et al.~\cite{zhang2022vehiclererank} extended this idea specifically
for vehicle ReID by incorporating orientation-guided query expansion, showing
that viewpoint-consistent neighbour selection further improves re-ranking
for highly viewpoint-sensitive queries. Moral et al.~\cite{moral2023syntheticrerank}
also studied vehicle-specific post-processing, including gallery video-clip
information in the re-ranking step on CityFlow-ReID and VeRi-776.

Average Query Expansion (AQE)~\cite{arandjelovic2012aqe} replaces each query
feature with the mean of its top-$k$ retrieved gallery features, smoothing the
query manifold and consistently improving mAP at negligible computational cost.
Database-side Augmentation (DBA) applies the same expansion to gallery features
and can be combined with AQE; we evaluate AQE alone as it provides the best
mAP--Rank-1 trade-off at $k{=}3$ for our ViT-B stream.

Supervised Contrastive Learning (SupCon)~\cite{khosla2020supcon} extends
contrastive losses to the supervised setting by treating all same-class samples
as positives within each mini-batch, producing more uniformly distributed and
discriminative embedding spaces than triplet loss alone~\cite{ye2021reidsurvey}.
We evaluate SupCon as an alternative metric objective for Stream~1 (variant
A5$\beta$); the deployed checkpoint (A5$\alpha$) retains a triplet + center
objective, and the two are shown to be statistically equivalent
(Section~\ref{sec:results}).

\subsection{Positioning}
To the best of the authors' knowledge, no prior work reports above 90\% mAP
on VeRi-776 using a pure ViT-B backbone---the best prior pure-ViT result
stands at 84.5\% mAP~\cite{li2023clipreid,opencodepapers2024veri}. The
present ViT-B recipe alone surpasses this by $+5.47$\,pp. No prior
two-stream fusion method has surpassed 93.3\% mAP on VeRi-776; the best
prior single-stream result stands at 92.9\% mAP~\cite{lu2025clipsenet}.

% =============================================================================
\section{Methodology}
\label{sec:method}
% =============================================================================

\subsection{Problem Formulation}
Given a query image $q$ and a gallery set
$\mathcal{G} = \{g_1, g_2, \ldots, g_N\}$, the task is to retrieve all
gallery images depicting the same vehicle identity as the query and rank them
above non-matching images. Performance is measured by mAP and Rank-1 CMC.
The standard VeRi-776 evaluation protocol~\cite{liu2018provid} employs 1\,678
query images and 11\,579 gallery images across 200 test identities; training
uses 37\,781 images of 576 identities. All reported results strictly follow
this protocol.

\subsection{Stream~1 --- Improved ViT-B Training Recipe}

\subsubsection{Backbone Architecture}
The backbone follows the TransReID architecture~\cite{he2021transreid}: a
ViT-Base/16 augmented with the SIE module and the JPM. The SIE module
encodes VeRi-776's 20 camera indices as learnable additive offsets to the
patch token sequence. The JPM operates as a parallel branch that randomly
shuffles groups of patch tokens and applies a separate classification head.
A BNNeck is applied on the 768-dimensional global CLS token before the
classification head~\cite{he2021transreid}.

The key deviation from the original TransReID recipe is the backbone
initialization: \texttt{vit\_base\_patch16\_clip\_224.openai} (OpenAI CLIP
pretrained ViT-B/16 via \texttt{timm}~\cite{radford2021clip}) replaces
ImageNet-21k-only weights. CLIP pretraining provides richer visual-semantic
representations that transfer more effectively to fine-grained vehicle
appearance.

\subsubsection{Loss Function}
Training employs a composite loss. The classification loss is cross-entropy
with label smoothing ($\varepsilon = 0.1$):
\begin{equation}
  \mathcal{L}_{\mathrm{CE\text{-}LS}} =
    -\sum_{c=1}^{C} y_c^{\mathrm{smooth}} \log \hat{p}_c, \;\;
    y_c^{\mathrm{smooth}} = (1{-}\varepsilon)y_c + \varepsilon/C
  \label{eq:cels}
\end{equation}
where $C = 576$ training identities. The \emph{deployed} headline checkpoint
(A5$\alpha$) pairs this classification loss with a soft-margin triplet loss
(margin $m{=}0.3$) and a delayed center loss
($\lambda{=}5\times10^{-4}$, enabled after epoch~30), under layer-wise
learning-rate decay 0.75:
\begin{equation}
  \mathcal{L}_{\mathrm{A5\alpha}} =
    \mathcal{L}_{\mathrm{CE\text{-}LS}}
    + \mathcal{L}_{\mathrm{tri}}(m{=}0.3)
    + \lambda\,\mathcal{L}_{\mathrm{center}}
  \label{eq:a5alpha}
\end{equation}
As a statistically-equivalent alternative, we also evaluate a
supervised-contrastive variant (A5$\beta$) that replaces the triplet and
center terms with Supervised Contrastive Loss
(SupCon)~\cite{khosla2020supcon} ($\tau = 0.07$, LLRD 0.65):
\begin{equation}
  \mathcal{L}_{\mathrm{SupCon}} =
    \sum_{i \in \mathcal{B}}
      \frac{-1}{|\mathcal{P}(i)|}
      \sum_{p \in \mathcal{P}(i)}
        \log
        \frac{\exp(\mathbf{z}_i \cdot \mathbf{z}_p / \tau)}
             {\sum_{a \in \mathcal{A}(i)} \exp(\mathbf{z}_i \cdot \mathbf{z}_a / \tau)}
  \label{eq:supcon}
\end{equation}
where $\mathcal{P}(i)$ is the set of positive pairs and $\mathcal{A}(i)$ all
other mini-batch samples, giving
$\mathcal{L}_{\mathrm{A5\beta}} = \mathcal{L}_{\mathrm{CE\text{-}LS}} + \mathcal{L}_{\mathrm{SupCon}}$.
A5$\alpha$ and A5$\beta$ are evaluated as separate ablation rows. The Wave~3
ablation campaign measured A5$\alpha$ at 89.06\% mAP and A5$\beta$ at 88.57\%
mAP under stream-1-only AQE+rerank evaluation; the 0.49\,pp difference
(A5$\alpha$ higher) falls within the $\pm 0.55$\,pp three-seed standard
deviation of A5$\alpha$ (88.75 $\pm$ 0.55\,pp mAP; see
Table~\ref{tab:ablation} and Section~\ref{sec:seed_variance}). Because the
difference falls well inside this three-seed standard deviation, the two
recipes are statistically equivalent and A5$\alpha$---the actually deployed
checkpoint---is retained as the headline Stream~1 recipe.
CircleLoss was observed to produce numerical instability (diverging training
loss) on CLIP-pretrained ViT-B backbones in our experiments and is therefore
excluded from the final recipe.

\subsubsection{Augmentation Pipeline}
Table~\ref{tab:augmentation} lists the full augmentation sequence applied
during Stream~1 training. All reported Stream~1 results use a
224\,$\times$\,224 input; hue/colour-corrupting augmentations (random
grayscale, strong hue jitter) were excluded as they remove vehicle colour
identity cues. A controlled 256\,$\times$\,256 \emph{training}-resolution
comparison was not reproduced in this campaign (see
Section~\ref{sec:discussion}).

\begin{table}[!t]
  \centering
  \caption{Stream~1 Training Augmentation Pipeline}
  \label{tab:augmentation}
  \begin{tabular}{clL{4.0cm}}
    \toprule
    Step & Transform & Parameters \\
    \midrule
    1 & Resize             & $224\times224$, bicubic \\
    2 & Random H-flip      & $p=0.5$ \\
    3 & Pad + Rand.\ crop  & padding=10\,px \\
    4 & ColorJitter        & bright.=0.2, cont.=0.15, sat.=0.1, hue=0.0 \\
    5 & Random Gauss.\ blur& kernel=5, $p=0.2$ \\
    6 & Random perspective & distortion=0.1, $p=0.2$ \\
    7 & Normalise          & CLIP $\mu$=(0.481,0.458,0.408),
                             $\sigma$=(0.269,0.261,0.276)~\cite{radford2021clip} \\
    8 & Random erasing     & $p=0.5$, scale=[0.02,0.33] \\
    \bottomrule
  \end{tabular}
\end{table}

\subsubsection{Training Hyperparameters}
Table~\ref{tab:hyperparams} compares the proposed Stream~1 configuration with
the published TransReID recipe~\cite{he2021transreid}. Training required
approximately 2.5--3.0 GPU-hours on $2\times$ NVIDIA T4 GPUs with
DataParallel. The exact $P{\times}K$ tuple is not recorded in the training
logs; only the effective batch size of 96 is confirmed.

\begin{table}[!t]
  \centering
  \caption{Training Hyperparameters: Stream~1 vs.\ TransReID~\cite{he2021transreid}}
  \label{tab:hyperparams}
  \begin{tabular}{L{2.2cm}L{2.3cm}L{2.5cm}}
    \toprule
    Hyperparameter & TransReID & Ours (Stream~1) \\
    \midrule
    Backbone        & ViT-B/16+SIE+JPM  & ViT-B/16+SIE+JPM \\
    Pre-train       & ImageNet-21k$\!\to\!$1K & OpenAI CLIP~\cite{radford2021clip} \\
    Input size      & $256\times256$    & $224\times224$ \\
    Optimizer       & SGD (mom=0.9)     & AdamW~\cite{loshchilov2019adamw} \\
    LR backbone     & 0.008             & $3.5\times10^{-4}$ (LLRD=0.75) \\
    LR head         & 0.008             & $3.5\times10^{-3}$ \\
    LR schedule     & Cosine decay      & Cosine + 10-ep warmup \\
    Batch size      & 64 (16$\times$4)  & 96 ($2\times$T4 DP) \\
    Epochs          & 120               & 140 \\
    Loss            & CE + Triplet      & CE-LS + Triplet + Center \\
    Mixed prec.     & fp16              & fp16 \\
    SIE cameras     & 20                & 20 \\
    JPM groups      & 4                 & 4 \\
    BNNeck          & Yes               & Yes \\
    \bottomrule
  \end{tabular}
\end{table}

\subsection{Stream~2 --- CLIP-SENet Branch}
The CLIP-SENet architecture~\cite{lu2025clipsenet} provides a complementary
feature stream. Where the ViT-B processes images as sequences of global patch
tokens through self-attention, CLIP-SENet processes images through a
ResNet101-IBN-a backbone~\cite{pan2018ibn} that applies local convolutional
filters with Instance-Batch Normalisation (IBN), capturing fine-grained local
texture patterns less salient in global attention maps.

The CLIP-SENet architecture consists of three components~\cite{lu2025clipsenet}:
\begin{enumerate}
  \item \textbf{CNN Backbone:} ResNet101-IBN-a~\cite{pan2018ibn} for
    2048-dimensional local appearance features. IBN balances instance
    normalisation (style invariance) and batch normalisation (discriminative
    learning), benefiting cross-camera ReID where camera-specific styles
    constitute nuisance variation.
  \item \textbf{Semantic Extraction Module (SEM):} The TinyCLIP image encoder
    extracts a 512-dimensional visual-semantic representation; unlike
    CLIP-ReID~\cite{li2023clipreid}, no text encoder is used at any
    stage~\cite{lu2025clipsenet}.
  \item \textbf{Adaptive Fine-Grained Enhancement Module (AFEM):} Adaptively
    groups and weights fine-grained vehicle attributes (model, colour, shape)
    from TinyCLIP, producing robust semantic embeddings without text
    annotations~\cite{lu2025clipsenet}.
\end{enumerate}

The CLIP-SENet v6 checkpoint is trained on VeRi-776 with Adam
($\mathrm{lr}\!=\!5\!\times\!10^{-4}$, wd$=\!5\!\times\!10^{-4}$, cosine
schedule), $P\!=\!8$, $K\!=\!8$, effective batch 128 (accum\_steps=2), CE-LS
+ SupCon, 24 epochs with 5-epoch warmup on a P100 GPU ($\approx$4\,h\,26\,min).
Input size is $320\!\times\!320$, matching the released CLIP-SENet~v6
checkpoint; a $256\!\times\!256$ \emph{training}-resolution variant was not
reproduced in this campaign.

\subsection{Post-processing Pipeline}
Post-processing is applied independently to each stream before fusion.

\subsubsection{Average Query Expansion (AQE)~\cite{arandjelovic2012aqe}}
For each query $\mathbf{q}$, the top-$k$ gallery features by cosine similarity
are averaged with the original query:
\begin{equation}
  \mathbf{q}_{\mathrm{AQE}} =
    \frac{1}{k{+}1}\!\left(\mathbf{q} + \sum_{i=1}^{k} \mathbf{g}_{\sigma(i)}\right)
  \label{eq:aqe}
\end{equation}
Stream~1 uses $k{=}3$; Stream~2 uses $k{=}10$, both selected by grid search
over $k\in\{1,3,5,10,20\}$. Stream~1's compact CLS token space is optimally
expanded with a tighter neighbourhood ($k{=}3$); Stream~2's broader
ResNet-IBN clusters benefit from a larger average ($k{=}10$), consistent with
the IBN design rationale for cross-camera style robustness~\cite{pan2018ibn}.

\subsubsection{$k$-Reciprocal Re-ranking~\cite{zhong2017rerank}}
The final distance for each query-gallery pair is:
\begin{equation}
  d_{\mathrm{final}}(q,g) =
    (1{-}\lambda)\,d_{\mathrm{Jaccard}}(q,g)
    + \lambda\,d_{\mathrm{cosine}}(q,g)
  \label{eq:rerank}
\end{equation}
Stream~1 (mAP-optimal): $k_1{=}80$, $k_2{=}15$, $\lambda{=}0.2$.
Stream~1 (Rank-1-optimal): $k_1{=}24$, $k_2{=}8$, $\lambda{=}0.2$.
Stream~2: $k_1{=}50$, $k_2{=}10$, $\lambda{=}0.1$. Parameters were selected
by exhaustive grid search over
$k_1\in\{20,22,24,25,26,28,30,50,80\}$,
$k_2\in\{6,7,8,9,10,15\}$,
$\lambda\in\{0.15,0.18,0.2,0.22,0.25,0.3\}$.

\subsubsection{Concat-patch + Flip TTA (Stream~1)}
Features are extracted under original and horizontally flipped inputs and
concatenated with patch mean pooling (CLS~+~mean patch tokens
$\to$~1536D). This configuration yields the best-mAP operating point
(89.97\% mAP and 97.79\% Rank-1), while the single-flip CLS-only 768D
configuration yields the best-Rank-1 operating point (85.14\% mAP and
98.33\% Rank-1).

\subsection{Score-Level Fusion}
Given post-processed similarity matrices $\mathbf{S}_{\mathrm{ViT}}$ and
$\mathbf{S}_{\mathrm{CLIP}}$, the fused similarity is:
\begin{equation}
  \mathbf{S}_{\mathrm{fusion}} =
    w_1\,\mathbf{S}_{\mathrm{ViT}} + w_2\,\mathbf{S}_{\mathrm{CLIP}},
  \quad w_1+w_2=1
  \label{eq:fusion}
\end{equation}
Concretely, the fusion pipeline uses the \textbf{1536D} concat-patch
flip-TTA Stream~1 representation (the best-mAP Stream~1 configuration of
Section~\ref{sec:method}, \emph{not} the 768D CLS-only variant) together
with the 2048D Stream~2 representation: AQE $k{=}3$ is applied to
\emph{each} stream, the resulting cosine-similarity matrices are combined
via Eq.~\eqref{eq:fusion}, and a \emph{single} joint $k$-reciprocal
re-ranking ($k_1{=}80, k_2{=}15, \lambda{=}0.2$) is applied to the fused
matrix. The fusion therefore re-applies a common AQE $k{=}3$ to both
streams; it does \emph{not} reuse Stream~2's standalone AQE $k{=}10$ /
re-ranking $50/10/0.1$, which characterise the standalone CLIP-SENet
result (91.44\% mAP) only.

Score-level fusion is preferred over feature-level fusion because the two
streams have fundamentally different feature geometries (1536D concat-patch
ViT representation vs.\ 2048D convolutional map) and normalisation schemes; combining at the
feature level would require a learned projection prone to overfitting on
576 training identities. The optimal weights $w_1{=}0.35$, $w_2{=}0.65$
were found by grid search on the VeRi-776 query/gallery split
($\approx$49.5\,min on T4). Since no separate held-out validation split
is used for weight selection, the chosen weights may reflect a mild
optimistic bias. However, Table~\ref{tab:fusion_sweep} demonstrates
that mAP remains above 93.0\% across $w_1 \in [0.2,\,0.5]$, indicating
that the optimal weight is not a sharp peak and that the practical impact
of this selection procedure on the reported result is negligible.

% =============================================================================
\section{Experimental Setup}
\label{sec:setup}
% =============================================================================

\subsection{Dataset}
VeRi-776~\cite{liu2016veri,liu2018provid} contains over 50\,000 images of
776 vehicles captured by 20 surveillance cameras covering a
$1.0\,\text{km}^{2}$ area over 24 hours. Each vehicle is annotated with
identity, camera index, viewpoint (8 categories), vehicle type, colour, and
spatiotemporal metadata~\cite{liu2018provid}. The standard split is:
37\,781 training images of 576 identities, 11\,579 gallery images of 200
test identities, and 1\,678 query images.

\subsection{Evaluation Protocol}
Results are reported under two conditions:
\begin{itemize}
  \item \textbf{w/o PP}: Cosine similarity, no re-ranking or query expansion.
  \item \textbf{w/ PP}: Full post-processing pipeline (AQE + $k$-reciprocal
    re-ranking + concat-patch flip TTA).
\end{itemize}
Primary metrics are mAP (\%) and Rank-1 (\%). All comparisons use matching
post-processing conditions.

\subsection{Implementation Details}
Table~\ref{tab:implementation} summarises the implementation configuration.

\begin{table}[!t]
  \centering
  \caption{Implementation Details}
  \label{tab:implementation}
  \begin{tabular}{L{3.2cm}L{4.5cm}}
    \toprule
    Component & Value \\
    \midrule
    Framework            & PyTorch 2.4.1+cu124, CUDA 12.4 \\
    S1 backbone          & \texttt{vit\_base\_patch16\_clip\_224.openai} \\
    S2 backbone          & ResNet101-IBN-a + TinyCLIP SEM \\
    S1 hardware          & $2\times$ NVIDIA T4, $\approx$2.5--3.0\,GPU-h \\
    S2 hardware          & NVIDIA P100, $\approx$4\,h\,26\,min \\
    Fusion sweep HW      & T4 (CPU-bound), $\approx$49.5\,min \\
    S1 feat.\ dim        & 1536D concat-patch+flip (best-mAP \& fusion);
                           768D CLS-only (best-R1) \\
    S2 feat.\ dim        & 2048D \\
    AQE $k$, standalone (S1 / S2) & 3 / 10 \\
    Re-rank S1 standalone (mAP)   & $k_1{=}80$, $k_2{=}15$, $\lambda{=}0.2$ \\
    Re-rank S1 standalone (R1)    & $k_1{=}24$, $k_2{=}8$, $\lambda{=}0.2$ \\
    Re-rank S2 standalone         & $k_1{=}50$, $k_2{=}10$, $\lambda{=}0.1$ \\
    Fusion post-proc     & AQE $k{=}3$ both streams $\to$ joint rerank
                           $k_1{=}80$, $k_2{=}15$, $\lambda{=}0.2$ \\
    Fusion weights       & $w_1{=}0.35$, $w_2{=}0.65$ \\
    \bottomrule
  \end{tabular}
\end{table}

\subsection{Reproducibility}
\label{sec:reproducibility}
Every number in this paper is produced by \emph{deterministic} evaluation
from two frozen, SHA-256-pinned checkpoints. Because inference---feature
extraction, AQE, $k$-reciprocal re-ranking, and score fusion---is
deterministic, each reported value reproduces \emph{exactly} on
re-execution. Table~\ref{tab:repro} lists the two released checkpoints, their
roles, and SHA-256 hashes.

\begin{table}[!t]
  \centering
  \caption{Released checkpoints and provenance. Both hashes are verified at
    evaluation time; the headline fusion (93.32\% mAP / 98.45\% Rank-1) is the
    deterministic evaluation of exactly these two files at $w_1{=}0.35$.}
  \label{tab:repro}
  \footnotesize
  \begin{tabular}{L{1.4cm}L{6.1cm}}
    \toprule
    Stream & Checkpoint, role, and SHA-256 \\
    \midrule
    Stream~1 & \texttt{vehicle\_transreid\_vit\_base\_veri776.pth} \\
             & ViT-B/16 CLIP; 89.97\% standalone (1536D), used at 1536D in fusion \\
             & {\scriptsize\texttt{8d32334a41933cbf0514fbc72e462124}} \\
             & {\scriptsize\texttt{f8785a8057d02b475aca8abf65dbe5e9}} \\
    \midrule
    Stream~2 & \texttt{clipsenet\_v6\_veri776.pth} \\
             & ResNet101-IBN-a + TinyCLIP (CLIP-SENet v6); 91.44\% standalone \\
             & {\scriptsize\texttt{d24bd3cd42cc4e90c7da71e4b955d7b7}} \\
             & {\scriptsize\texttt{e3274614e53fc686153b02c116d20a1c}} \\
    \bottomrule
  \end{tabular}
\end{table}

We distinguish two reproducibility regimes. (i)~\textbf{Exact reproduction:}
evaluating the released checkpoints with the released notebooks reproduces
every table value bit-for-bit, since inference is deterministic and the
checkpoints are hash-pinned---this is the operative guarantee for all headline
numbers, and the fusion headline was confirmed identical across three
independent evaluations (two GPU notebooks and a CPU recomputation from dumped
features). (ii)~\textbf{Within-variance reproduction:} retraining a stream
from scratch reproduces its recipe only within the $\pm0.55$\,pp three-seed
standard deviation (Section~\ref{sec:seed_variance}), because GPU training is
not bit-reproducible (cuDNN autotuning, non-deterministic atomics) even under
fixed seeds; the released checkpoint---not a fresh retrain---therefore defines
the headline. The standard VeRi-776 protocol is asserted at evaluation time
($1\,678$ query, $11\,579$ gallery, $576$ training identities); all evaluation
used PyTorch~2.4.1 (CUDA~12.4) on a single NVIDIA~T4 GPU. The evaluation
notebooks, the checkpoint manifest (with full hashes), and the per-table
provenance mapping are released in the project repository
(\url{https://github.com/MRKDaGods/gp}).

% =============================================================================
\section{Results and Analysis}
\label{sec:results}
% =============================================================================

\subsection{Comparison with State of the Art}

Table~\ref{tab:sota} compares the proposed single-stream (S1) and fusion
results against published methods on VeRi-776.

\begin{table}[!t]
  \centering
  \caption{State-of-the-Art Comparison on VeRi-776.
    \cmark\,=\,post-processing applied; \xmark\,=\,cosine similarity only.}
  \label{tab:sota}
  \begin{tabular}{L{2.4cm}L{1.6cm}L{3.2cm}C{1.3cm}C{1.2cm}C{1.0cm}L{1.0cm}}
    \toprule
    Method & Venue & Backbone & mAP~(\%) & R1~(\%) & w/~PP & Ref. \\
    \midrule
    PROVID       & TMM 2018   & CNN (multi-modal)         & 53.42 & 81.56 & \xmark & \cite{liu2018provid} \\
    VAMI         & CVPR 2018  & CNN                       & 50.13 & 77.03 & \xmark & \cite{zhou2018vami} \\
    PVEN         & CVPR 2020  & ResNet50                  & 79.50 & 95.60 & \xmark & \cite{meng2020pven} \\
    HPGN         & T-ITS 2022 & CNN+Graph                 & 80.18 & 96.72 & \xmark & \cite{shen2022hpgn} \\
    VehicleNet   & TMM 2020   & ResNet50                  & 83.41 & 96.78 & \cmark & \cite{zheng2020vehiclenet} \\
    TransReID    & ICCV 2021  & DeiT-B/16                 & 82.30 & 97.10 & \xmark & \cite{he2021transreid} \\
    CLIP-ReID    & AAAI 2023  & ViT-B/16                  & 84.50 & 97.30 & \xmark & \cite{li2023clipreid} \\
    RPTM         & WACV 2023  & ResNet                    & 88.00 & 97.30 & \xmark & \cite{ghosh2023rptm} \\
    MBR4B-LAI    & ITSC 2023  & ResNet50 (multi-branch)   & 92.09 & 98.03 & \cmark & \cite{luo2023mbr4b} \\
    CLIP-SENet   & arXiv 2025 & ResNet101-IBN-a+TinyCLIP  & 92.90 & 98.70 & \cmark & \cite{lu2025clipsenet} \\
    \midrule
    \textbf{Ours (S1, w/ PP)} & ---  & ViT-B/16 (CLIP init)              & \textbf{89.97} & \textbf{97.79} & \cmark & --- \\
    \textbf{Ours (Fusion)}    & ---  & ViT-B + ResNet101-IBN-a + TinyCLIP & \textbf{93.32} & \textbf{98.45} & \cmark & --- \\
    \bottomrule
  \end{tabular}
\end{table}

The proposed ViT-B recipe surpasses all prior pure-ViT and transformer-only
methods by a substantial margin. The fusion result establishes a new state of
the art in mAP on VeRi-776, with $+0.42$\,pp mAP over CLIP-SENet~\cite{lu2025clipsenet}
and $+1.23$\,pp mAP over MBR4B-LAI~\cite{luo2023mbr4b}.

\subsection{Pure ViT-B Comparison}

Table~\ref{tab:vit_compare} isolates the comparison against prior
transformer-based methods using the same ViT-B/16 backbone family.

\begin{table}[!t]
  \centering
  \caption{Transformer-Based Methods on VeRi-776}
  \label{tab:vit_compare}
  \begin{tabular}{L{2.4cm}L{1.4cm}L{1.9cm}C{1.0cm}C{1.0cm}}
    \toprule
    Method & Venue & Backbone & mAP & R1 \\
    \midrule
    TransReID (ViT-B)~\cite{he2021transreid}   & ICCV'21 & ViT-B/16  & 80.6 & 97.0 \\
    TransReID* (DeiT-B)~\cite{he2021transreid} & ICCV'21 & DeiT-B/16 & 82.3 & 97.1 \\
    CLIP-ReID~\cite{li2023clipreid}             & AAAI'23 & ViT-B/16  & 84.5 & 97.3 \\
    \midrule
    \textbf{Ours (S1, w/ PP)}                  & ---     & \textbf{ViT-B/16 (CLIP)} & \textbf{89.97} & \textbf{97.79} \\
    \bottomrule
  \end{tabular}
\end{table}

The proposed recipe delivers $+5.47$\,pp mAP over
CLIP-ReID~\cite{li2023clipreid} and $+7.67$\,pp over TransReID*.
For clarity, we report best-mAP and best-R1 configurations separately: the
89.97\% mAP / 97.79\% Rank-1 result corresponds to a concat-patch flip TTA
configuration with AQE $k{=}3$ and $k$-reciprocal re-ranking $(k_1{=}80, k_2{=}15,
\lambda{=}0.2)$, whereas the 98.33\% Rank-1 / 85.14\% mAP result comes from
a single-flip CLS-only configuration with $k$-reciprocal re-ranking only
$(k_1{=}24, k_2{=}8, \lambda{=}0.2)$ and no AQE.

\subsection{Ablation Study}

Table~\ref{tab:ablation} traces the complete performance trajectory. Endpoint
values are verified against published results~\cite{he2021transreid} and
repository experiment logs; all intermediate rows are measured from isolated
ablation experiments.

\begin{table}[!t]
  \centering
  \caption{Stepwise Ablation: 82.3\% $\to$ 93.3\% mAP on VeRi-776.
    Rows A1--A5$\beta$ report isolated-component contributions measured
    under a common operating point (Stream~1 only, AQE $k{=}3$ +
    $k$-reciprocal re-ranking, single-flip 768D CLS); deltas are relative
    to baseline A1. The cumulative trajectory below combines components
    progressively and adds concat-patch flip TTA, CLIP-SENet, and fusion.%
    \protect\footnotemark}
  \label{tab:ablation}
  \begin{tabular}{L{4.2cm}C{1.1cm}C{1.0cm}C{1.2cm}}
    \toprule
    Configuration & mAP & R1 & $\Delta$ vs A1 \\
    \midrule
    \multicolumn{4}{l}{\emph{Isolated Component Ablation (Stream~1, AQE+rerank)}} \\
    A1 --- Baseline: DeiT + SGD + CE + Triplet            & 57.62 & 80.45 & --- \\
    A2 --- +CLIP init only                                & 70.97 & 86.83 & $+13.35$ \\
    A3 --- +SupCon only (DeiT + SGD + CE-LS + SupCon)    & 59.48 & 79.74 & $+1.86$ \\
    A4 --- +AdamW/LLRD only (DeiT + CE + Triplet)        & 80.91 & 94.82 & $+23.29$ \\
    A5$\alpha$ --- Full (CLIP + AdamW/LLRD=0.75 + CE-LS + Triplet + CenterLoss; \emph{deployed}) & 89.06 & 97.74 & $+31.44$ \\
    A5$\beta$ --- Full (CLIP + AdamW/LLRD=0.65 + CE-LS + SupCon; alternative) & 88.57 & 96.66 & $+30.95$ \\
    \midrule
    \multicolumn{4}{l}{\emph{Cumulative Trajectory $\to$ Headline}} \\
    TransReID baseline (DeiT-B*)~\cite{he2021transreid} & 82.30 & 97.10 & --- \\
    Base cosine, ViT-B (no PP)                           & 82.22 & 97.50 & $-0.08$ \\
    \quad+AQE $k{=}3$                                    & 84.08 & 97.20 & $+1.86$ \\
    \quad+Re-rank $k_1{=}80,k_2{=}15,\lambda{=}0.2$     & 89.63 & 97.74 & $+5.55$ \\
    \quad+Concat-patch + flip TTA (1536D; A5$\alpha$, deployed) & 89.97 & 97.79 & $+0.34$ \\
    \midrule
    CLIP-SENet v6, base cosine~\cite{lu2025clipsenet}    & 82.34 & 96.54 & --- \\
    \quad+AQE $k{=}10$                                   & 89.21 & 96.90 & $+6.87$ \\
    \quad+Re-rank $k_1{=}50,k_2{=}10,\lambda{=}0.1$     & 91.44 & 97.08 & $+2.23$ \\
    \midrule
    \textbf{Fusion ($w_1{=}0.35$, $w_2{=}0.65$)}        & \textbf{93.32} & \textbf{98.45} & $+1.88$ \\
    \bottomrule
  \end{tabular}
\end{table}
\footnotetext{All reported numbers are reproduced from frozen,
  SHA-256-pinned checkpoints via deterministic evaluation notebooks;
  isolated-ablation rows (A1--A5$\beta$) are from-scratch retrains within the
  $\pm 0.55$\,pp three-seed band. See Section~\ref{sec:reproducibility} for
  checkpoint hashes and the two reproducibility regimes.}

The isolated-component block clarifies the recipe drivers. AdamW with
layer-wise LR decay (A4) is the single largest training-recipe contributor,
adding $+23.29$\,pp mAP over the SGD baseline A1; CLIP initialisation (A2)
is the second largest at $+13.35$\,pp. SupCon in isolation (A3) barely
moves the needle ($+1.86$\,pp) but combines well once CLIP and AdamW are in
place: the deployed A5$\alpha$ recipe (CLIP + AdamW + CE-LS + Triplet +
CenterLoss) reaches 89.06\% mAP, statistically indistinguishable from the
alternative A5$\beta$ recipe (CLIP + AdamW + CE-LS + SupCon,
88.57\% mAP)---a 0.49\,pp single-run difference that lies within the
$\pm 0.55$\,pp three-seed standard deviation of A5$\alpha$
(88.75\,$\pm$\,0.55\,pp; Section~\ref{sec:seed_variance}).
Within the cumulative trajectory, the most impactful post-processing step
is $k$-reciprocal re-ranking~\cite{zhong2017rerank} on Stream~1
($+5.55$\,pp); AQE~\cite{arandjelovic2012aqe} adds $+1.86$\,pp; TTA
contributes $+0.34$\,pp. Fusion provides $+1.88$\,pp above the stronger
post-processed CLIP-SENet stream, confirming genuine complementarity
between the two feature spaces.

\subsection{Training Seed Variance}
\label{sec:seed_variance}

To quantify training-time stochasticity, the A5$\alpha$ recipe was retrained
from scratch with three additional seeds $\{42, 123, 456\}$ under identical
hyperparameters, each evaluated at the isolated-ablation operating point
(Stream~1 only, single-flip 768D, AQE $k{=}3$ + $k$-reciprocal re-ranking;
Table~\ref{tab:ablation}, top block). The purpose of Table~\ref{tab:seed_variance}
is to bound the recipe's \emph{training} variance and so underpin the
A5$\alpha$/A5$\beta$ equivalence argument; it is not a re-derivation of the
released headline checkpoint, whose reproducibility is established
separately (below).

\begin{table}[!t]
  \centering
  \caption{A5$\alpha$ Training Seed Variance on VeRi-776 (Stream~1 only,
    single-flip 768D, AQE $k{=}3$ + $k$-reciprocal re-ranking). Quantifies
    from-scratch retraining noise of the recipe.}
  \label{tab:seed_variance}
  \begin{tabular}{lcc}
    \toprule
    Seed & mAP~(\%) & R1~(\%) \\
    \midrule
    42  & 89.48 & 97.26 \\
    123 & 88.14 & 97.08 \\
    456 & 88.64 & 97.20 \\
    \midrule
    \textbf{Mean $\pm$ std (n=3)} & \textbf{88.75 $\pm$ 0.55} & \textbf{97.18 $\pm$ 0.07} \\
    \bottomrule
  \end{tabular}
\end{table}

The three-seed standard deviation of 0.55\,pp mAP exceeds the 0.49\,pp
A5$\alpha$--A5$\beta$ single-run gap, confirming that the two recipes are
statistically equivalent under realistic training noise. This is purely a
training-time analysis: the headline Stream~1 result (89.97\% mAP) is
\emph{not} the mean of these retrains but the deterministic evaluation of the
released, SHA-pinned checkpoint at the best-mAP operating point (1536D
concat-patch flip TTA). Its reproducibility is therefore exact---the same
frozen checkpoint yields the same number on every run, independent of
training variance---which is the guarantee we provide by releasing the
checkpoint. For transparency, that checkpoint scores 89.63\% at the 768D
operating point of this table, at the favourable end of the retrain
distribution; a reproducer training from scratch should expect the band
above, whereas a reproducer evaluating the released checkpoint obtains the
headline exactly.

\subsection{Fusion Weight Sensitivity}

\begin{table}[!t]
  \centering
  \caption{Fusion Weight Sensitivity on VeRi-776 (1536D Stream~1,
    no-TTA Stream~2)}
  \label{tab:fusion_sweep}
  \begin{tabular}{cccc}
    \toprule
    $w_1$ (ViT-B) & $w_2$ (CLIP-SENet) & mAP~(\%) & R1~(\%) \\
    \midrule
    0.20 & 0.80 & 93.13 & 98.03 \\
    0.25 & 0.75 & 93.25 & 98.15 \\
    0.30 & 0.70 & 93.21 & 98.15 \\
    \textbf{0.35} & \textbf{0.65} & \textbf{93.32} & \textbf{98.45} \\
    0.40 & 0.60 & 93.30 & 98.39 \\
    0.45 & 0.55 & 93.26 & 98.45 \\
    0.50 & 0.50 & 93.10 & 98.45 \\
    \bottomrule
  \end{tabular}
\end{table}

mAP remains above 93.0\% for $w_1 \in [0.2, 0.5]$, demonstrating robustness
to the exact weight choice. The consistent superiority of all fusion
configurations over either stream in isolation confirms that the two streams
capture complementary information about vehicle identity.

We additionally verified that the Stream~1 feature dimensionality has
negligible effect on the fused result. The 1536D concat-patch flip-TTA
representation and the 768D single-flip CLS-only representation fuse to
within 0.01\,pp mAP at their respective optima (1536D: 93.32\% at
$w_1{=}0.35$; 768D: 93.31\% at $w_1{=}0.30$), even though the 1536D
representation is 0.34\,pp stronger standalone (89.97\% vs.\ 89.63\%). The
standalone TTA gain is thus largely redundant with the complementary
CLIP-SENet signal; we retain the 1536D representation as it is both the
best-mAP standalone Stream~1 configuration and marginally best in fusion.

\subsection{Post-processing Contribution Analysis}

Table~\ref{tab:pp_contrib} breaks down the per-stream contribution of each
post-processing step.

\begin{table}[!t]
  \centering
  \caption{Per-Stream Post-processing Contribution}
  \label{tab:pp_contrib}
  \begin{tabular}{L{3.4cm}L{1.6cm}C{1.0cm}C{1.0cm}}
    \toprule
    Configuration & Stream & mAP & R1 \\
    \midrule
    Base cosine (no PP)              & ViT-B      & 82.22 & 97.50 \\
    +AQE $k{=}3$                     & ViT-B      & 84.08 & 97.20 \\
    +Re-rank                         & ViT-B      & 89.63 & 97.74 \\
    +TTA (1536D)                     & ViT-B      & 89.97 & 97.79 \\
    \midrule
    Base cosine (no PP)              & CLIP-SENet & 82.34 & 96.54 \\
    +AQE $k{=}10$                    & CLIP-SENet & 89.21 & 96.90 \\
    +Re-rank                         & CLIP-SENet & 91.44 & 97.08 \\
    \midrule
    \textbf{Full PP $\to$ Fusion}    & Both       & \textbf{93.32} & \textbf{98.45} \\
    \bottomrule
  \end{tabular}
\end{table}

Both streams begin at similar base performance ($\approx$82\% mAP);
the larger relative gain of CLIP-SENet from AQE ($+6.87$\,pp vs.\ $+1.86$\,pp
for ViT-B) reflects the broader nearest-neighbour clusters of the
ResNet-IBN feature space, which benefit more from graph-based smoothing.

\subsection{Qualitative Analysis}

Visual inspection indicates the fusion model handles frontal-to-rear viewpoint
transitions, illumination shifts, and partial occlusion well. Failure cases
arise under: (a)~near-identical vehicles of the same make, model, and colour
at similar viewpoints, where both streams return similar-appearing but
incorrect identities; (b)~extreme occlusion ($>$60\% vehicle body);
and (c)~very low-resolution gallery crops. In these cases, both streams tend
to fail on the same queries, indicating correlated rather than complementary
errors in the hardest cases. Table~\ref{tab:qual} summarizes these failure
modes alongside representative success cases.

% ── Qualitative error-mode table ──
\begin{table}[!t]
  \caption{Qualitative Failure Mode Analysis on VeRi-776 (200 Test Identities)}
  \label{tab:qual}
  \centering
  \small
  \begin{tabular}{p{2.0cm}p{2.6cm}p{2.6cm}}
    \toprule
    \textbf{Failure Mode} & \textbf{ViT-B stream} & \textbf{CLIP-SENet stream} \\
    \midrule
    Same make/model/colour, identical viewpoint &
      Returns visually similar but incorrect identity &
      Returns same incorrect identity; errors correlated \\[2pt]
    Extreme occlusion ($>$60\% body) &
      Rank-1 miss; partial hood/trunk features dominate &
      Rank-1 miss; AFEM misses occluded part regions \\[2pt]
    Very low-resolution gallery ($<$40$\times$40 px) &
      Patch tokens degenerate; CLS unstable &
      IBN-a loses texture; both streams fail jointly \\[2pt]
    Frontal-to-rear transition (success case) &
      SIE camera offset reduces cross-view gap & 
      IBN style normalisation handles camera shift \\[2pt]
    Illumination shift day/night (success case) &
      CLIP semantic prior abstracts away lighting &
      SupCon anchor pulls same-ID clusters together \\
    \bottomrule
  \end{tabular}
  \vspace{2pt}
  {\footnotesize Success cases are included for contrast.
   Failure modes (rows 1--3) show correlated errors across both streams,
   indicating a ceiling reachable only by improved cross-view invariance.
   Full retrieval panel figures are available in the project
   repository.}
\end{table}

% =============================================================================
\section{Discussion}
\label{sec:discussion}
% =============================================================================

\subsection{What Drives the ViT-B Gain}
The $+7.67$\,pp improvement over the published TransReID baseline arises from
independent and additive components. CLIP pretraining initialisation is the
primary driver, providing a rich visual-semantic prior that enables better
fine-grained vehicle discrimination from the outset, consistent with transfer
learning observations in the ReID
literature~\cite{li2023clipreid,radford2021clip}. For context, initialising
with DINOv2~\cite{oquab2023dinov2}---a self-supervised alternative---was
observed to underperform the CLIP-pretrained ViT-B in our preliminary
experiments, suggesting that CLIP's \emph{semantic supervision}---not merely
pretraining scale---is a key factor. The metric loss (triplet + center),
AQE, and $k$-reciprocal re-ranking contribute further additive gains as
detailed in the ablation study.

\subsection{Why Score-Level Fusion Works}
The complementarity is architectural. The ViT-B backbone processes images as
global patch-token sequences with multi-head self-attention, producing
representations sensitive to holistic vehicle structure and long-range spatial
relationships~\cite{he2021transreid,dosovitskiy2021vit}. The CLIP-SENet
ResNet101-IBN-a backbone~\cite{pan2018ibn} applies hierarchical local
convolutional filters capturing fine-grained texture, with IBN providing
simultaneous style normalisation and discriminative training. The AFEM branch
injects attribute-level structure from TinyCLIP~\cite{lu2025clipsenet}
orthogonal to ViT patch attention. Score-level fusion preserves each stream's
post-processing calibration independently and requires no additional training.

The optimal $w_2{=}0.65$ for CLIP-SENet reflects its higher standalone
post-processed performance (91.44\% vs.\ 89.97\% mAP for ViT-B). The ViT-B
branch still contributes $+1.88$\,pp above CLIP-SENet alone, confirming
genuine complementarity.

\subsection{Backbone Scale and Overfitting}
ViT-L/14 (304M parameters) was evaluated as a higher-capacity alternative.
Despite its larger parameter count, it was observed to overfit VeRi-776's
576 training identities, falling below the ViT-B CLIP result after re-ranking
in our experiments. This motivates systematic post-processing over backbone
scaling as the primary performance lever for small-scale ReID benchmarks.

\subsection{Limitations}
\begin{enumerate}
  \item The full post-processing pipeline (AQE + re-ranking) has
    $\mathcal{O}(N^{2})$ complexity in gallery size $N$ and is primarily suited
    to offline evaluation rather than real-time deployment.
  \item The fusion strategy requires two inference passes, approximately
    doubling memory and latency relative to a single-stream model.
  \item All results are reported on VeRi-776 under the standard supervised
    protocol; cross-dataset generalisation to VeRi-Wild or CityFlowV2 is not
    directly achieved. In particular, a direct port of the fusion mechanism to
    CityFlowV2 produced monotonic MTMC IDF1 degradation, with a best-case
    change of approximately $-1.8$\,pp at higher fusion weights. This suggests
    that the VeRi-776--CityFlowV2 domain gap remains the binding constraint and
    that improved cross-camera invariance, rather than additional fusion, is
    required for MTMC.
  \item An earlier draft of the Stream~1 training recipe described the metric
    loss and layer-wise learning-rate decay imprecisely. The present version
    corrects that disclosure by reporting isolated component ablations for
    CLIP initialisation (A2), SupCon versus triplet loss (A3), and AdamW with
    LLRD (A4), together with dual Full references: the actually deployed
    CE-LS + Triplet + CenterLoss recipe (A5$\alpha$, 89.06\% mAP) and the
    alternative CE-LS + SupCon recipe (A5$\beta$, 88.57\% mAP). The
    0.49\,pp gap between them lies within the $\pm 0.55$\,pp three-seed
    standard deviation of A5$\alpha$ (88.75 $\pm$ 0.55\,pp mAP), so the two
    recipes are statistically equivalent and A5$\alpha$---the deployed
    checkpoint---is retained as the headline Stream~1 recipe. The two-stream
    fusion cost and the VeRi-776--CityFlowV2 domain-gap constraint remain
    as noted above.
\end{enumerate}

% =============================================================================
\section{Conclusion}
\label{sec:conclusion}
% =============================================================================

This paper presented two complementary contributions to vehicle
re-identification on VeRi-776. A systematically improved ViT-Base/16
training recipe---combining CLIP pretraining initialisation, a tuned
metric-learning objective, an optimised augmentation pipeline, AdamW
optimisation, and a three-step post-processing pipeline---achieves 89.97\%
mAP and 97.79\% Rank-1 at its best-mAP operating point, with a separate
best-Rank-1 operating point of 85.14\% mAP and 98.33\% Rank-1. This is the
best reported mAP for any pure ViT-B model on this benchmark, surpassing
TransReID by $+7.67$\,pp mAP~\cite{he2021transreid}. Score-level fusion with
a CLIP-SENet branch~\cite{lu2025clipsenet} achieves 93.32\% mAP and 98.45\%
Rank-1, a new state of the art in mAP on VeRi-776, exceeding the
previous best published mAP (CLIP-SENet standalone: 92.9\%~\cite{lu2025clipsenet})
by $+0.42$\,pp. Rank-1 accuracy (98.45\%) is competitive with but does not
exceed the CLIP-SENet standalone Rank-1 (98.7\%).

The comprehensive ablation from 82.3\% to 93.3\% mAP demonstrates that
systematic recipe refinement---rather than architectural novelty---yields
substantial performance gains on established benchmarks. All code and
checkpoints are publicly released.

Future directions include: cross-dataset generalisation on
VeRi-Wild~\cite{lou2019veriwild}; knowledge distillation from the fusion model
into a single-stream model; and extension to nighttime surveillance and
CityFlowV2.

\bibliographystyle{IEEEtran}
\bibliography{references}

\end{document}
